\chapter{Oh Ho}

After the war, drugs were everywhere.

They surfaced the way ruins did—half-buried, scattered, unclaimed. Weapons once engineered to confuse enemy minds now lay dormant in drawers, tins, and coat pockets. Some were lethal. Others were merely strange. Peter Sands had learned to tell the difference.

He collected them carefully.

Not to escape reality, as others assumed, but to press against it. He believed—quietly, stubbornly—that some of the so-called hallucinations were not false at all. They were glimpses. Cracks in a membrane that normally sealed human perception shut.

Green pills. Red capsules. White tablets scored into halves or quarters. He arranged them like an apothecary, hiding them from everyone but himself. And he experimented.

He favored the frightening visions. Just barely. He wanted to look over the edge without falling. Fear, to him, was proof that something real lay beyond.

He remembered his father once gripping the handles of a shock machine at an amusement park, inching them apart to increase the current. A contest with pain. Watching his father endure it had filled Pete with awe. His father had let go eventually—but only after proving something.

Pete knew he wasn’t like that. But he wanted to be.

So he mixed drugs the way an alchemist once mixed metals. Carefully. With safeguards. Always with another person present—Lurine—so she could stop him if he went too far. He knew the maps of the brain. Knew which glands to quiet, which regions to release. He wanted experience, not destruction.

“I’m nuts,” he once told her.

“Then stop,” she said.

He couldn’t. He was searching.

For what, he couldn’t fully articulate. Something that lay behind a curtain. Something that death itself might tear open—but he wanted to see it before that.

“You could build instruments,” Lurine told him. “Measure it.”

“I don’t want data,” he said. “I want contact.”

She worried he would go too far one day and never come back.

“I can’t wait,” he told her. “It doesn’t come this side of the grave.”

He despised the Servants of Wrath for calling survival proof of righteousness. If their god was evil, then surely he spared the wicked. By their own logic, they were damned by being alive.

That bored Pete. So he turned back to his pills.

One night, Lurine asked what he was really seeking.

He hesitated. Then told her about the sting.

Not poison. Not pain in the modern sense. But penetration. A piercing awareness.

Once, under the influence, he had become violent—had broken things, hurt her—and then it had happened. Something had entered him like a barbed hook, a gaff driven through his side. The pain was absolute. Unmistakably real.

Above him, he had seen them.

Three figures. Calm. Watching. Not cruel, not kind. Simply there.

They held the spear until he woke.

The purpose of the pain had been awakening.

Ever since, he feared it—and sought it.

That night, he swallowed a new combination.

And he saw a man.

A powerful, young figure with white hair, dressed like something out of antiquity. The man smiled warmly and opened a massive book filled with unreadable symbols. Meaningless to Pete, though heavy with importance.

Then the vision vanished.

A voice replaced it.

Something simpler.

A clay pot floated before him—plain, unglazed, humble.

“I’ll tell you my name,” it said. “I’m Oh Ho.”

The pot explained patiently: the grand figure had been a composite intelligence, ancient, impressive, deliberately obscure. It impressed through mystery.

But Oh Ho identified itself. Treated Pete as an equal.

“That which is benign names itself,” Oh Ho said. “That which will not, is not.”

Before leaving, it offered a final message: Sophia will return.

The drugs wore off.

Pete sat in his living room again. Lurine smoked quietly. Nothing had changed.

“I saw Him,” Pete said.

She nodded without interest.

Later, she told him the Servants of Wrath were sending Tibor McMasters on a pilgrimage—to find the God of Wrath and capture his likeness for a mural.

“You should sabotage it,” she said calmly. “Break the cow’s leg. Drain the battery.”

Pete scoffed—until the doorbell rang.

Dr. Abernathy stood outside.

The Christian priest spoke gently. Offered instruction. Confirmation. Rituals that required time, patience, humility.

Poker followed. Old coins. Quiet talk.

Tibor arrived later, trembling with fear. He confessed that if he joined the Christian Church, he might not have to go on the pilgrimage.

Abernathy looked at him steadily.

“Fear often masks guilt,” he said.

Tibor had no guilt. Only terror.

Confession was offered anyway—not as escape, but as confrontation.

The idea unsettled Tibor. And yet, somehow, eased him.

When Abernathy invited him to come privately the next morning, Tibor agreed. Something unseen felt lighter.

The game resumed.

Silver coins changed hands.

Outside, the cow-cart waited.

And the pilgrimage—whether toward truth or ruin—drew closer.

